\section{The later earthquake comparison}
\label{sec:earthquake-comparison}

The analytical theorem concerns intrinsic product densities of an ideal stationary process. The numerical study in Section~\ref{sec:simulation} concerns the consistency of a sampler and estimators with that theorem. Neither task by itself compares the model with the finite-catalogue ratios in the earthquake studies. Such a comparison must reproduce the observation window, event selection, location uncertainty, tuple bins and random-reference construction before its ratios can be interpreted.

The required historical observation mask and complete event-selection rule have not yet been recovered. In particular, the illustrative cube of side $378$ km used in earlier count experiments is not an established footprint for the reported $803$-event catalogue. We therefore do not put those counts into the numerical validation study or treat their difference as a fitted-model rejection. The matched earthquake experiment remains planned rather than completed.

When the observation design is available, the comparison should keep three quantities separate: the intrinsic correlation $g_k$, the ratio of expected observed and reference tuple counts, and the expectation of a catalogue-normalised ratio. They agree only under additional conditions. The same study should report count dispersion, empty-window frequency, nearest-neighbour behaviour and shape distributions, since the reciprocal identities constrain only selected factorial correlations. Empty and low-count catalogues must remain in the sample, and one fixed model parameter set must be used at every order.

Even before that comparison, the expanded witness shows an important limitation. Its intensity is approximately $1.5713\times10^{-9}\,\mathrm{km}^{-3}$, so a cube of side $378$ km has expected count about $0.0849$. This calculation describes the witness in an illustrative window; it is not a reconstruction of the historical experiment. A useful earthquake model would also need temporal triggering, magnitude marks, fault geometry and an observation process, none of which is supplied by the present spatial existence construction.
